% Generated draft: requires mathematical and source review.
\documentclass[11pt]{article}
\usepackage[margin=25mm]{geometry}
\usepackage{fontspec}
\setmainfont{Latin Modern Roman}
\usepackage{amsmath,amssymb,mathtools,unicode-math}
\setmathfont{Latin Modern Math}
\usepackage{enumitem,needspace,xurl,hyperref}
\setlength{\parindent}{0pt}
\setlength{\parskip}{5pt}
\setlength{\emergencystretch}{3em}
\newcommand{\N}{\mathbb{N}}
\newcommand{\Z}{\mathbb{Z}}
\newcommand{\Q}{\mathbb{Q}}
\newcommand{\R}{\mathbb{R}}
\newcommand{\C}{\mathbb{C}}
\newcommand{\F}{\mathbb{F}}
\newcommand{\A}{\mathbb{A}}
\newcommand{\PP}{\mathbb{P}}
\newcommand{\E}{\mathbb{E}}
\newcommand{\eps}{\varepsilon}
\newcommand{\dd}{\,\mathrm{d}}
\begin{document}
\section*{UnsolvedMath ID: 800243}
\textbf{Source problem code:} catalogue entry
\subsection{Definitions and mathematical statement}
Let \(n\geq 1\) be an integer and let \(\Omega\subset\C^n\) be a domain, meaning a nonempty connected open set. Call \(\Omega\) pseudoconvex if it admits a continuous plurisubharmonic exhaustion function. Here plurisubharmonic means that the restriction to each complex line is subharmonic, and exhaustion means that every sublevel set is relatively compact in \(\Omega\). Say that \(\partial\Omega\) is smooth if, near every boundary point, there is a real-valued \(C^\infty\) defining function \(\rho\) such that \(\Omega=\{\rho<0\}\) locally and \(\dd\rho\neq0\) on the boundary.

Under the infinitesimal interpretation of the metrics, put \(\Delta=\{\zeta\in\C:|\zeta|<1\}\), and let \(\mathcal O(X,Y)\) denote the holomorphic maps from \(X\) to \(Y\). For \(z\in\Omega\) and \(v\in\C^n\), define the Carath\'eodory and Kobayashi infinitesimal pseudometrics by
\[
 c_\Omega(z;v)=\sup_{f\in\mathcal O(\Omega,\Delta)}\frac{|\dd f_z(v)|}{1-|f(z)|^2},
 \qquad
 k_\Omega(z;v)=\inf\bigl\{|\lambda|:\lambda\in\C,\ g\in\mathcal O(\Delta,\Omega),\ g(0)=z,\ \lambda g'(0)=v\bigr\}.
\]
They are globally comparable if there exists \(C\geq1\) such that
\[
 C^{-1}k_\Omega(z;v)\leq c_\Omega(z;v)\leq Ck_\Omega(z;v)
 \quad\text{for every }z\in\Omega\text{ and }v\in\C^n.
\]
Since \(c_\Omega\leq k_\Omega\), failure of comparability is equivalently
\[
 \forall C\geq1\ \exists z\in\Omega\ \exists v\in\C^n\setminus\{0\}:\quad
 k_\Omega(z;v)>C\,c_\Omega(z;v).
\]

\textbf{Problem, under this interpretation.} Does there exist a positive integer \(n\) and a pseudoconvex domain \(\Omega\subset\C^n\) with smooth boundary for which the displayed failure of global comparability holds?

\textbf{Unresolved source qualifications.} The source says ``smoothly bounded'' without explicitly stating whether \(\Omega\) must also be bounded. If boundedness is intended, add that hypothesis to the existence question. The source also does not specify whether the metrics are infinitesimal pseudometrics or their associated distance functions, nor whether \(n\) is fixed or may be chosen. The displayed formulation uses infinitesimal pseudometrics, a single comparison constant valid throughout the domain and in every tangent direction, and existential quantification over \(n\); these interpretations require review against the intended source meaning.
\subsection{Short English statement}
The question attributed to S. Krantz asks for a smoothly bounded pseudoconvex domain where the Kobayashi and Carath\'eodory metrics fail to be comparable. Under the stated infinitesimal interpretation, no single constant bounds the Kobayashi metric by the Carath\'eodory metric everywhere and in every tangent direction. The intended dimension, boundedness requirement, and meaning of metric remain to be confirmed.
\subsection{Sources}
\begin{itemize}
\item Local UnsolvedMath registry, permanent ID \(800243\), source problem code \texttt{catalogue entry}, title: \emph{Several Complex Variables Problem List --- Problem 41}. The supplied record labels the problem open; this label is not an independently verified mathematical or literature assessment.
\item Buzzard, Fornaess, and Krantz, \emph{The Several Complex Variables Problem List} (1995), Problem 41. Supplied source URL: \url{https://arxiv.org/abs/math/9509202}. The record reports access on 2026-10-16 and extraction by \texttt{arxiv\_latex}; these are supplied provenance details, not a claim of independent source access or verification.
\item The extracted question is signed S. Krantz, with supplied contact \texttt{SK@math.wustl.edu}. The registry lists Buzzard, Fornaess, and Krantz in its proposer field and gives proposed year 1995; the signature specifically attributes this question to S. Krantz.
\end{itemize}
\end{document}
